\scanpage{9}

\mspage{4}

Les str. de Brown (\unsure{cofibration catégorie|cofibration catégories|cofibration-catégorie} de Brown) permettent de pousser l'alg. homotopique très loin, semble-t-il, en ce qui concerne les structures \deleted{\unsure{impliquées|exactes}} qui tournent autour des suites \unsure{exactes|coexactes}, de suspension, et elles ont \unsure{la|l'air} \unsure{vertu|même} d'être stables par passage de $M$ à $S\backslash M$. Par contre, elles ne sont absolument pas prévues pour la stabilité par passage aux $M/S$, ni pour les suites exactes de cosuspension. \deleted{Mais} Si on \deleted{\unsure{devient}} \unsure{le|la|les} demande, alors l'axiome de factorisation \inserted{b)} $f=pi$ $(\unsure{i|j}\in TC,p\in F)$ est exigé dans \deleted{\unsure{les}} \unsure{sa|se} force (la forme a) restant faible, $f=pi$ avec $i\in C$ et $p\in W$ seulement, pas forcément $p\in F$ et \deleted{\unsure{on}} \unsure{a} fortiori pas \deleted{\unsure{son}} exiger $p\in TF$). Cela ne suffit pas, \deleted{que} faute d'exiger $F\subset\mathrm{Fib}_W$, pour assurer l'existence d'une bonne théorie de fibres $W$-homotopique (alors que l'on en a une, par les seules \unsure{cofibration-catégories|cofibrations-catégories}, pour les \uline{co}fibres $W$-homotopiques).

% archive sheet number: 8

